\chapter{Introduction}

This blueprint organizes the Lean~4 formalization of the main theorem of
Deng--Shen, \emph{The four-dimensional Anderson model: a case study for
critical SPDEs} (arXiv:2607.10105; all references use the v1 numbering).
The accompanying documentation provides complementary views.  The
\href{https://github.com/sairmath/anderson4d_lean/blob/main/docs/DESIGN.md}
{\texttt{DESIGN.md}} describes architecture and scope;
\href{https://github.com/sairmath/anderson4d_lean/blob/main/docs/PAPER_MAP.md}
{\texttt{PAPER\_MAP.md}} gives node IDs and paper-level dependencies (the
\texttt{uses}-labels here reuse those IDs verbatim);
\href{https://github.com/sairmath/anderson4d_lean/blob/main/docs/PAPER_TO_LEAN.md}
{\texttt{PAPER\_TO\_LEAN.md}} gives exact paper-number-to-public-declaration
links; and
\href{https://github.com/sairmath/anderson4d_lean/blob/main/docs/PAPER_INDEX.md}
{\texttt{PAPER\_INDEX.md}} is the generated module inventory.  This chapter
states the target and the global conventions; it contains no formalizable nodes.

\section{The target}

Let \(\xi\) be spatial white noise on \(\mathbb{T}^4 = [-\pi,\pi]^4\), let
\(\rho\) be a smooth cutoff (Section~\ref{sec:conventions}), and
\(\xi_\varepsilon = \rho_\varepsilon * \xi\). Consider
\[
\mathcal{L}_\varepsilon \;=\; 1 - \Delta - \lambda_\varepsilon\,\xi_\varepsilon
  + \mathcal{C}_\varepsilon,
\qquad \lambda_\varepsilon = \frac{\lambda}{\sqrt{\lvert\log\varepsilon\rvert}},
\]
with the explicit renormalization constants
\(\mathcal{C}_\varepsilon = \sum_{q \le A}\mathcal{C}_{2q}\),
\(A = \lfloor\lvert\log\varepsilon\rvert\rfloor\), of
Chapter~\ref{ch:detparametrix}. The paper's Theorem~1.1 asserts that
\(\mathcal{H}_\varepsilon = \lambda_\varepsilon^{-1}(G_\varepsilon - G)\) converges in
law, as \(\varepsilon \to 0^+\), to an explicit centered Gaussian field.

\textbf{Formal main statement} (design in \texttt{DESIGN.md}~\S2): for every
cutoff \(\rho\) there is \(\lambda_0(\rho) > 0\) such that for
\(0 < \lambda < \lambda_0\) and every finite family of Fourier modes
\((\alpha_j,\beta_j)_{j\le s}\), the random vectors
\(\big(\widehat{\mathcal{H}_\varepsilon}(\alpha_j,\beta_j)\big)_{j\le s}\)
converge in distribution to the corresponding Gaussian limit
(Chapter~\ref{ch:main}). This finite-mode form is exactly what the paper's
proof establishes at (3.34); no distribution-space topology enters.

\textbf{External input.} The single statement quoted by the paper from
Gabriel--Rosati (arXiv:2602.22509) --- Proposition~3.6, our node
\ref{P-3.6} --- is treated as an explicit hypothesis: it is stated in Lean
in full, marked \emph{external} here, never assumed as an axiom, and the
main theorem \texttt{main\_conditional} takes it as an argument. No
unconditional \texttt{main} is declared; a proof of the cited
Gabriel--Rosati result is outside the scope of this repository.

\section{Conventions}\label{sec:conventions}

The following conventions are fixed once and used silently everywhere;
getting them wrong invalidates constants globally, so they form the fixed
``normalization ledger''.

\begin{enumerate}
\item \textbf{Torus and measure.} \(\mathbb{T}^4\) is the product of four
  copies of \(\mathbb{R}/2\pi\mathbb{Z}\); the reference measure is Lebesgue
  measure of total mass \((2\pi)^4\). mathlib's \texttt{AddCircle} Fourier
  API is normalized against probability Haar measure; the translation
  dictionary (all \(2\pi\)-factors) is recorded with the Lean definitions.
\item \textbf{Characters and coefficients.} \(e_k(x) = e^{i k\cdot x}\);
  for kernels, \(\widehat K(\alpha,\beta)
  = \iint e^{i(\alpha\cdot x + \beta\cdot y)}K(x,y)\,\mathrm{d}x\,
  \mathrm{d}y = \langle e_{-\alpha},\, K e_\beta\rangle\) with the inner
  product conjugate-linear in the first slot (paper (3.23); sign frozen by
  a rank-one roundtrip test).
\item \textbf{Cutoff.} \(\rho\colon\mathbb{R}^4 \to \mathbb{R}\) smooth,
  nonnegative, supported in the ball of an \emph{explicit} radius \(R\)
  (no unit-ball normalization), invariant under coordinate permutations
  and sign flips (the class \(\mathcal{E}\), node \ref{D-E}), with
  \(\int\rho = 1\). Torus objects arise by \(2\pi\mathbb{Z}^4\)-periodization;
  mollification is equivalently the Fourier multiplier
  \(\widehat{\rho}(\varepsilon k)\).
\item \textbf{White noise.} Defined through Fourier coefficients: two
  independent standard real Gaussians per \(\{\pm k\}\)-orbit and one for
  the zero mode, assembled so that
  \(\mathbb{E}[g_k g_l] = \mathbf{1}_{k=-l}\) and
  \(\mathbb{E}[g_k \overline{g_l}] = \mathbf{1}_{k=l}\). Since these
  characters are orthonormal for probability Haar while the paper uses
  Lebesgue measure of mass \((2\pi)^4\), the physical field is
  \(\xi_\varepsilon=(2\pi)^{-2}\sum_k
  \widehat\rho(\varepsilon k)g_ke_k\). This half-density is part of the
  normalization ledger.
\item \textbf{Junk values.} Partial operations are totalized by junk
  values (\(G_\varepsilon := 0\) off the invertibility event, junk limit law
  for \(\lambda^2 \ge 2\pi^2\), etc.); every estimate carries the
  hypotheses that exclude the junk branch.
\item \textbf{Constants.} No asymptotic notation in Lean statements.
  Constants referenced by several nodes are named definitions
  (\(C_0 > 1000\), \(D = e^{C_0^{10}}\), \dots); one-off line constants
  are existentially quantified. All constants may depend on \(\rho\)
  (including its support radius \(R\)).
\item \textbf{Word model.} All permutation sums of the paper's \S5.4 are
  stated over \emph{words} (maps with prescribed fiber sizes). The
  definitional ledger
  \(\mathrm{paperSum} = \big(\prod_{\mathfrak{l}} m_{\mathfrak{l}}!\big)
  \cdot \mathrm{wordSum}\) (node \ref{D-ledger}) aligns statements with
  the paper's labeled normalization: paper-facing statements put
  \(\mathrm{paperSum}\) on the left with the paper's right-hand sides
  verbatim.
\end{enumerate}
